\cvsection{Профиль}
Senior Java Backend Engineer с опытом более 15 лет в телекоммуникациях, стриминговых
сервисах и retail-системах. Развиваю бизнес-критичные backend-системы без остановки
продуктового развития. Зона ответственности охватывает бизнес-логику и интеграции,
модернизацию runtime и данных, а также надежность production-систем. Принимаю
технические решения и отвечаю за их реализацию в production.

\cvsection{Ключевые компетенции}
\begin{itemize}[leftmargin=1.3em, itemsep=0.2em]
  \item Backend: Java 8--21, Spring Boot, REST, SOAP, gRPC
  \item Системы: распределенные развертывания, интеграции, авторизация, API gateway, модернизация legacy
  \item Данные: PostgreSQL, Oracle, оптимизация SQL, JPA/Hibernate, Liquibase
  \item Поставка: GitLab CI/CD, Docker, Docker Swarm, Linux, Nginx
  \item Наблюдаемость: Prometheus, Grafana, Micrometer
\end{itemize}

\cvsection{Опыт работы}
\cvrole{Senior Java Developer / Team Lead}{2016 -- настоящее время}{Рестрим Медиа (Wink)}{Москва, Россия}
\begin{itemize}[leftmargin=1.3em, itemsep=0.2em]
  \item Отвечал за backend подписок, покупок и интеграций Wink с аудиторией около 10 млн пользователей в месяц; каждый макрорегион обрабатывал около 300 RPS и хранил около 2 ТБ данных.
  \item Реализовал новые способы покупки и подписки, развил интеграционные потоки и выделил B2B-функции в отдельный подпроект.
  \item Спроектировал и создал сервисы авторизации и API gateway для основного продукта.
  \item Вел поэтапную модернизацию с Java 6 до Java 21 и Spring Boot 3.5 без остановки продуктового развития.
  \item Заменил разрозненные Oracle-скрипты управляемыми Liquibase-миграциями с проверками, rollback, release tags, процедурами и триггерами; оптимизировал SQL и переработал DAO-слой.
  \item Развил событийное взаимодействие через NSQ: типизированные сообщения, дедупликация, transactional hooks и persisted outbox.
  \item Добавил Actuator, метрики Prometheus/Micrometer, readiness, трассировку запросов и интеграционные тесты с Oracle Testcontainers.
  \item Определял технические решения для команды из пяти инженеров, провел около 20 собеседований и нанял пять инженеров.
\end{itemize}

\cvrole{Senior Java Developer}{2013 -- 2016}{AT Consulting}{Москва, Россия}
\begin{itemize}[leftmargin=1.3em, itemsep=0.2em]
  \item Разработка движка расчета скидок для retail-продуктов и сервисных API.
  \item Разработка сервисов расчета KPI для telecom-систем.
  \item Разработка функциональности систем управления знаниями для контакт-центров.
\end{itemize}

\cvrole{Опыт разработки}{2008 -- 2013}{Icicall / Intellectika / AuditNT}{Москва, Россия}

\cvsection{Образование}
Диплом специалиста по прикладной математике и информатике, 5-летняя программа\\
Южный федеральный университет (бывший РГУ), 2003 -- 2008

\cvsection{Языки}
Английский (рабочий), Русский (родной), Французский (начальный)
